\Chapter{Propozycje rozwiązań systemu Electronic Medical Record}\label{chapter:project-idea}

\setlength{\parskip}{0pt}
\setlength{\parindent}{1.25cm}

\section{Wprowadzenie do procesu projektowego}
Na podstawie przeprowadzonej analizy funkcjonowania Szpitala im. Bajgla przygotowano kilka możliwych kierunków rozwiązania problemów związanych z organizacją dokumentacji medycznej. Ze względu na skalę placówki oraz złożoność procesów medycznych nie istnieje jedno uniwersalne rozwiązanie, które w jednakowy sposób odpowiadałoby na wszystkie potrzeby organizacyjne, techniczne i finansowe szpitala.

Proces projektowy został więc oparty na analizie różnych wariantów wdrożenia systemu EMR (Electronic Medical Record), różniących się zakresem funkcjonalności oraz poziomem zaawansowania technologicznego. Przygotowane rozwiązania uwzględniają zarówno aktualne potrzeby szpitala, jak i możliwość dalszego rozwoju systemu w przyszłości.

Istotnym aspektem procesu projektowego było również uwzględnienie potrzeb różnych grup użytkowników, w tym lekarzy, pielęgniarek, laboratoriów, administracji oraz samych pacjentów.

Przygotowane warianty obejmują trzy poziomy wdrożenia systemu: rozwiązanie podstawowe koncentrujące się na cyfryzacji najważniejszych procesów, rozwiązanie optymalne stanowiące pełny zintegrowany system EMR dla dużego szpitala oraz rozwiązanie rozszerzone obejmujące dodatkowe funkcje automatyzacji i inteligentnego zarządzania placówką.


\section{Założenia projektowe}
Każdy z analizowanych wariantów systemu został opracowany na podstawie założeń projektowych oraz wymagań związanych z funkcjonowaniem nowoczesnych systemów ochrony zdrowia. Niezależnie od zakresu wdrożenia wszystkie rozwiązania muszą zapewniać możliwość centralnego przechowywania dokumentacji medycznej oraz obsługi wielu oddziałów w ramach jednego środowiska informatycznego.

Ważnym założeniem projektowym jest również dostosowanie systemu do pracy różnych grup użytkowników. Interfejs powinien być intuicyjny oraz możliwie prosty w codziennym użytkowaniu, tak aby zmniejszyć ryzyko błędów wynikających z nieprawidłowego korzystania. Projekt zakłada także możliwość korzystania z systemu na urządzeniach stacjonarnych lub mobilnych wykorzystywanych przez personel medyczny podczas pracy na oddziałach.

Każdy z wariantów musi również spełniać wymagania dotyczące bezpieczeństwa danych medycznych oraz zgodności z regulacjami obowiązującymi na terenie Unii Europejskiej. System powinien umożliwiać kontrolę dostępu do dokumentacji, rozdzielenie uprawnień użytkowników, monitorowanie aktywności w systemie oraz zabezpieczenie danych przed utratą lub nieautoryzowanym dostępem.
\section{Wariant I – rozwiązanie podstawowe}

\subsection{Charakterystyka rozwiązania}
Wariant podstawowy zakłada cyfryzację najważniejszych procesów związanych z prowadzeniem dokumentacji medycznej w Szpitalu im. Bajgla. Rozwiązanie koncentruje się przede wszystkim na odejściu od papierowego przechowywania danych pacjentów oraz stworzeniu centralnego miejsca umożliwiającego przeglądanie podstawowej dokumentacji medycznej przez personel szpitala.
System w tej wersji ma charakter uproszczony i obejmuje wyłącznie najważniejsze funkcje niezbędne do poprawy organizacji pracy placówki. 

Projekt zakłada wykorzystanie stanowisk komputerowych znajdujących się na oddziałach bez wdrażania rozbudowanej infrastruktury mobilnej. Wariant ten został przygotowany jako rozwiązanie wymagające stosunkowo niewielkich zmian organizacyjnych oraz ograniczonego zakresu wdrożenia.

Rozwiązanie podstawowe skupia się głównie na ograniczeniu problemów związanych z przechowywaniem dokumentacji papierowej oraz uproszczeniu dostępu do informacji o pacjentach. System nie obejmuje jednak zaawansowanej automatyzacji procesów ani pełnej integracji wszystkich elementów funkcjonowania szpitala.
\subsection{Zakres funkcjonalności}
Wariant podstawowy obejmuje centralną bazę danych pacjentów umożliwiającą przechowywanie podstawowej dokumentacji medycznej w formie elektronicznej. Personel medyczny posiada możliwość przeglądania historii leczenia pacjenta, wyników badań oraz podstawowych informacji dotyczących hospitalizacji z poziomu stanowisk komputerowych znajdujących się na oddziałach.

System umożliwia również elektroniczne dodawanie wyników badań laboratoryjnych oraz dokumentacji diagnostycznej. W ramach rozwiązania przewidziano podstawowy model logowania użytkowników oraz podział uprawnień pomiędzy personelem medycznym a administracją.

Wariant obejmuje również uproszczony portal pacjenta umożliwiający dostęp do wybranych wyników badań oraz podstawowych informacji dotyczących leczenia. 

\subsection{Zalety rozwiązania}
Największą zaletą wariantu podstawowego jest stosunkowo niski poziom złożoności wdrożenia oraz ograniczone wymagania infrastrukturalne. Rozwiązanie pozwala szybko ograniczyć część problemów wynikających z dokumentacji papierowej bez konieczności przeprowadzania całkowitej reorganizacji pracy szpitala.

Wariant podstawowy może również zmniejszyć opór części personelu wobec cyfryzacji, ponieważ zakres zmian organizacyjnych pozostaje stosunkowo niewielki. System jest prostszy w utrzymaniu i wymaga mniejszej liczby szkoleń niż bardziej rozbudowane rozwiązania.

Dodatkową zaletą jest możliwość stopniowego rozwijania systemu w przyszłości poprzez dodawanie kolejnych funkcji wraz z dalszą cyfryzacją placówki.

\subsection{Ograniczenia rozwiązania}
Wariant podstawowy nie eliminuje całkowicie problemów związanych z przepływem informacji pomiędzy oddziałami oraz ograniczoną mobilnością personelu medycznego. Brak wykorzystania tabletów może powodować dalsze opóźnienia w dostępie do danych w sytuacjach wymagających szybkiej reakcji.

Ograniczeniem rozwiązania pozostaje również niski poziom automatyzacji procesów administracyjnych oraz brak zaawansowanych funkcji wspierających organizację pracy dużej placówki medycznej. System nie zapewnia pełnej integracji wszystkich procesów szpitalnych, przez co część działań nadal wymagać będzie ręcznego przekazywania informacji pomiędzy użytkownikami.

Wariant podstawowy może również okazać się rozwiązaniem niewystarczającym w perspektywie dalszego rozwoju placówki oraz rosnących wymagań związanych z cyfryzacją ochrony zdrowia.

\section{Wariant II – rozwiązanie optymalne}

\subsection{Charakterystyka rozwiązania}
Wariant optymalny zakłada wdrożenie pełnego zintegrowanego systemu EMR obejmującego wszystkie kluczowe obszary funkcjonowania szpitala. Rozwiązanie zostało zaprojektowane z myślą o dużej wielooddziałowej placówce wymagającej sprawnego przepływu informacji pomiędzy personelem medycznym, laboratoriami, diagnostyką oraz administracją.

System w tej wersji ma stanowić centralne środowisko pracy dla całego szpitala, umożliwiające bieżący dostęp do dokumentacji pacjentów w czasie rzeczywistym. Projekt zakłada wykorzystanie zarówno stanowisk komputerowych, jak i urządzeń mobilnych wykorzystywanych przez personel medyczny podczas pracy na oddziałach.

Rozwiązanie obejmuje pełną cyfryzację dokumentacji medycznej oraz automatyzację wybranych procesów organizacyjnych związanych z obiegiem informacji i wyników badań.

\subsection{Zakres funkcjonalności}
Wariant optymalny obejmuje centralną bazę danych pacjentów, elektroniczną dokumentację medyczną, elektroniczne zlecenia badań oraz bieżącą synchronizację danych pomiędzy oddziałami. Personel medyczny posiada możliwość przeglądania i aktualizowania danych pacjentów w czasie rzeczywistym z poziomu komputerów oraz tabletów.

System umożliwia automatyczne przekazywanie wyników badań laboratoryjnych i diagnostycznych do dokumentacji pacjenta oraz monitorowanie statusu wykonywanych badań. W rozwiązaniu przewidziano również moduł powiadomień informujący personel o nowych wynikach, zmianach dokumentacji lub ważnych zdarzeniach związanych z leczeniem pacjenta.

Projekt obejmuje rozbudowany system ról i poziomów dostępu, historię zmian dokumentacji, backup danych oraz mechanizmy monitorowania aktywności użytkowników. Pacjenci posiadają dostęp do portalu umożliwiającego sprawdzanie wyników badań, historii leczenia oraz wybranych informacji dotyczących hospitalizacji.


\subsection{Zalety rozwiązania}
Największą zaletą wariantu optymalnego jest pełna integracja procesów związanych z dokumentacją medyczną i przepływem informacji pomiędzy oddziałami. System znacząco usprawnia komunikację pomiędzy personelem medycznym.

Rozwiązanie poprawia organizację pracy szpitala poprzez umożliwienie szybkiego dostępu do aktualnych danych pacjentów. Wykorzystanie urządzeń mobilnych zwiększa wygodę pracy personelu i umożliwia dostęp do dokumentacji bez konieczności powrotu do stacjonarnych stanowisk komputerowych.

Wariant optymalny zapewnia również wysoki poziom bezpieczeństwa danych medycznych oraz możliwość dalszej rozbudowy systemu w przyszłości. Rozwiązanie odpowiada potrzebom dużej placówki medycznej i pozostaje zgodne z kierunkiem rozwoju nowoczesnych systemów eHealth w Unii Europejskiej.

\subsection{Ograniczenia rozwiązania}
Wdrożenie wariantu optymalnego wymaga większych nakładów organizacyjnych oraz technicznych niż rozwiązanie podstawowe. Konieczne staje się przygotowanie odpowiedniej infrastruktury sieciowej, zakup urządzeń mobilnych oraz przeprowadzenie szkoleń dla personelu medycznego.

Rozwiązanie wiąże się również z większym poziomem zmian organizacyjnych w codziennym funkcjonowaniu placówki. Część pracowników może początkowo odczuwać trudności związane z przejściem na w pełni cyfrowy model pracy.

Wariant wymaga także bardziej rozbudowanego wsparcia technicznego oraz regularnego utrzymania systemu ze względu na większą liczbę funkcjonalności i użytkowników.

\section{Wariant III – rozwiązanie rozszerzone}

\subsection{Charakterystyka rozwiązania}
Wariant rozszerzony zakłada rozbudowę systemu EMR o dodatkowe funkcje analizy danych oraz częściowej automatyzacji organizacji pracy szpitala. W przeciwieństwie do wariantu optymalnego rozwiązanie nie skupia się wyłącznie na elektronicznej dokumentacji i komunikacji pomiędzy oddziałami, ale dodatkowo umożliwia bieżące monitorowanie działania placówki oraz analizowanie danych organizacyjnych.

System ma wspierać dyrekcję i administrację w kontrolowaniu funkcjonowania szpitala poprzez wykorzystanie bardziej zaawansowanych narzędzi raportowania, monitorowania i analizy danych generowanych podczas codziennej pracy placówki.
\subsection{Zakres funkcjonalności}
Wariant rozszerzony obejmuje dodatkowe moduły analityczne oraz funkcje automatyzacji procesów organizacyjnych. System może automatycznie tworzyć raporty dotyczące obciążenia oddziałów, liczby pacjentów, czasu oczekiwania na badania lub wykorzystania personelu medycznego.

Rozwiązanie może również generować inteligentne powiadomienia dotyczące problemów organizacyjnych, takich jak opóźnienia badań, przeciążenie oddziałów lub brakujące elementy dokumentacji. Dodatkowo przewidziano możliwość integracji z wybranymi urządzeniami medycznymi, dzięki czemu część danych mogłaby trafiać do systemu automatycznie bez konieczności ręcznego przepisywania wyników.

W projekcie uwzględniono także bardziej rozbudowaną aplikację pacjenta umożliwiającą szerszy kontakt ze szpitalem oraz dostęp do większej liczby informacji dotyczących leczenia.
\subsection{Zalety rozwiązania}
Największą zaletą wariantu rozszerzonego jest możliwość wykorzystania danych do wspomagania organizacji pracy szpitala oraz szybszego wykrywania problemów związanych z funkcjonowaniem placówki. System może ułatwić analizowanie obciążenia oddziałów oraz usprawnić zarządzanie personelem i zasobami szpitala.

Dodatkową zaletą jest możliwość dalszej rozbudowy systemu o kolejne funkcje automatyzacji oraz integracji z dodatkowymi urządzeniami i modułami w przyszłości.
\subsection{Ograniczenia rozwiązania}
Wariant rozszerzony wymaga bardziej rozbudowanej infrastruktury technicznej oraz większego wsparcia administracyjnego niż wariant optymalny. Rozwiązanie wiąże się również z większym poziomem skomplikowania utrzymania systemu oraz koniecznością bardziej zaawansowanego zarządzania jego działaniem.

Dodatkowo w przypadku integracji mechanizmów analizy danych opartych na sztucznej inteligencji zaistniałaby konieczność dostosowania systemu do wymagań wynikających z regulacji AI Act oraz przepisów dotyczących ochrony danych osobowych. Oznaczałoby to między innymi konieczność odpowiedniego dokumentowania sposobu działania systemu, monitorowania wykorzystania modeli AI oraz zapewnienia przejrzystości procesu przetwarzania danych medycznych.

Konieczne byłoby również informowanie pacjentów oraz pracowników o wykorzystaniu mechanizmów sztucznej inteligencji w analizie danych medycznych i organizacyjnych. W zależności od zakresu działania systemu mogłoby być wymagane także wdrożenie dodatkowych procedur związanych z nadzorem nad działaniem modeli AI oraz bezpieczeństwem danych.

\section{Porównanie wariantów}

\subsection{Tabela porównawcza}

\renewcommand{\arraystretch}{1.3}

\begin{longtable}{|p{4cm}|p{3cm}|p{3cm}|p{3cm}|}

\caption{Porównanie wariantów systemu EMR}
\label{tab:warianty_emr} \\

\hline
\textbf{Kryterium} &
\shortstack{\textbf{Wariant I} \\ podstawowy} &
\shortstack{\textbf{Wariant II} \\ optymalny} &
\shortstack{\textbf{Wariant III} \\ rozszerzony} \\
\hline
\endfirsthead

\hline
\textbf{Kryterium} &
\shortstack{\textbf{Wariant I} \\ podstawowy} &
\shortstack{\textbf{Wariant II} \\ optymalny} &
\shortstack{\textbf{Wariant III} \\ rozszerzony} \\
\hline
\endhead

Zakres cyfryzacji dokumentacji &
częściowy &
pełny &
pełny \\
\hline

Centralna baza danych pacjentów &
tak &
tak &
tak \\
\hline

Elektroniczna dokumentacja medyczna &
podstawowa &
pełna &
pełna \\
\hline

Integracja oddziałów &
ograniczona &
pełna &
pełna \\
\hline

Dostęp do danych w czasie rzeczywistym &
częściowy &
tak &
tak \\
\hline

Mobilny dostęp dla personelu &
brak &
tablety i urządzenia mobilne &
tablety i urządzenia mobilne \\
\hline

Portal pacjenta &
podstawowy &
rozbudowany &
rozbudowany \\
\hline

Automatyczne przekazywanie wyników badań &
częściowe &
tak &
tak \\
\hline

Moduły analityczne i raportowanie &
brak &
podstawowe raporty &
zaawansowana analiza danych \\
\hline

Analiza obciążenia oddziałów &
brak &
brak &
rozbudowana \\
\hline

Inteligentne powiadomienia i alerty &
brak &
podstawowe &
rozbudowane \\
\hline

Integracja z urządzeniami medycznymi &
brak &
ograniczona &
rozszerzona \\
\hline

Wsparcie organizacji pracy szpitala &
podstawowe &
wysokie &
bardzo wysokie \\
\hline

Bezpieczeństwo danych i kontrola dostępu &
podstawowe &
wysokie &
bardzo wysokie \\
\hline

Wykorzystanie analizy danych AI &
brak &
brak &
tak \\
\hline

Wymagania związane z AI Act &
nie dotyczy &
nie dotyczy &
wymagane \\
\hline

Poziom automatyzacji &
niski &
średni &
wysoki \\
\hline

Stopień skomplikowania wdrożenia &
niski &
wysoki &
bardzo wysoki \\
\hline

Wymagania infrastrukturalne &
niskie &
średnie &
wysokie \\
\hline

Zakres szkoleń personelu &
niewielki &
duży &
bardzo duży \\
\hline

\end{longtable}

\subsection{Analiza porównawcza}

Przeprowadzona analiza pokazuje, że każdy z wariantów odpowiada na potrzeby szpitala w innym stopniu. Wariant podstawowy umożliwia częściową cyfryzację dokumentacji oraz ograniczenie ilości dokumentów papierowych, jednak nie zapewnia pełnej integracji procesów medycznych i administracyjnych. Rozwiązanie to charakteryzuje się najniższymi kosztami wdrożenia i najmniejszym poziomem skomplikowania infrastruktury.

Wariant optymalny zapewnia pełną cyfryzację dokumentacji medycznej, integrację oddziałów oraz sprawny dostęp do danych pacjentów w czasie rzeczywistym. Rozwiązanie oferuje szeroki zakres funkcjonalności przy jednoczesnym zachowaniu umiarkowanych wymagań organizacyjnych i infrastrukturalnych.

Wariant rozszerzony oferuje najwyższy poziom automatyzacji oraz dodatkowe funkcje związane z analizą danych i integracją elementów sztucznej inteligencji. Rozwiązanie wiąże się jednak z największym poziomem złożoności wdrożenia, wysokimi kosztami utrzymania oraz dodatkowymi wymaganiami prawnymi związanymi między innymi z AI Act.

\section{Wybór rozwiązania docelowego}

Na podstawie przeprowadzonej analizy jako rozwiązanie docelowe wybrano wariant optymalny. Rozwiązanie to stanowi kompromis pomiędzy zakresem funkcjonalności, kosztami wdrożenia oraz możliwościami organizacyjnymi szpitala.

Wybrany wariant spełnia wszystkie kluczowe wymagania placówki związane z pełną cyfryzacją dokumentacji medycznej, poprawą przepływu informacji oraz zwiększeniem dostępności danych pacjentów. Jednocześnie nie wymaga tak rozbudowanej infrastruktury oraz tak dużych zmian organizacyjnych jak wariant rozszerzony.

Klient zdecydował się na rozwiązanie, które jest możliwe do wdrożenia w obecnych warunkach finansowych szpitala, a jednocześnie zapewnia możliwość dalszej rozbudowy systemu w przyszłości.
\section{Kryteria oceny rozwiązania}
Po wprowadzeniu docelowego wariantu systemu EMR zostanie przeprowadzona ocena skuteczności wdrożonego rozwiązania na podstawie określonych kryteriów funkcjonalnych, organizacyjnych oraz technicznych. Celem oceny będzie sprawdzenie, w jakim stopniu system usprawni funkcjonowanie szpitala oraz ograniczy problemy związane z papierowym obiegiem dokumentacji.

Pierwszym kryterium będzie poziom cyfryzacji dokumentacji medycznej. Zostanie on określony jako procent dokumentów obsługiwanych elektronicznie w stosunku do całej dokumentacji powstającej w codziennej pracy szpitala. Pod uwagę zostaną wzięte między innymi karty pacjentów, wyniki badań, zlecenia lekarskie, raporty pielęgniarskie, wypisy oraz dokumenty administracyjne. Im większy udział dokumentów obsługiwanych w systemie cyfrowym, tym wyższy będzie poziom cyfryzacji placówki.

Drugim kryterium będzie czas dostępu do dokumentacji pacjenta. Pomiar zostanie wykonany poprzez sprawdzenie, ile czasu potrzebuje uprawniony użytkownik na odnalezienie historii leczenia, aktualnych wyników badań lub ostatnich zleceń lekarskich. Analiza będzie obejmowała typowe sytuacje pracy szpitala, takie jak przyjęcie pacjenta, obchód lekarski lub przygotowanie wypisu.

Trzecim kryterium będzie dostępność systemu dla użytkowników w codziennej pracy. Ocenie zostanie poddana liczba stanowisk oraz urządzeń umożliwiających korzystanie z systemu, możliwość pracy mobilnej oraz dostępność danych na oddziałach, na których personel nie pracuje stale przy komputerze. Szczególnie istotne będzie sprawdzenie, czy lekarz lub pielęgniarka będą mogli uzyskać dostęp do danych pacjenta bez konieczności opuszczania oddziału lub przerywania wykonywanych czynności.

Czwartym kryterium będzie bezpieczeństwo i kontrola dostępu do danych. W ramach oceny sprawdzone zostanie, czy system posiada role użytkowników, indywidualne konta, historię zmian dokumentacji, rejestr logowań oraz możliwość ograniczania dostępu do danych zgodnie z zakresem obowiązków pracownika. Kryterium to będzie mierzone między innymi liczbą dostępnych poziomów uprawnień oraz możliwością monitorowania aktywności użytkowników w systemie.
